
Multilingual pretraining datasets routinely apply CCNet-style perplexity quality filters, but the magnitude of the language-group retention disparity produced by applying a single global threshold to pre-computed ccnet\_perplexity scores on RedPajama-V2 has not been precisely quantified using a standard association metric. We find that this common practice produces a language-family retention divide rather than a language-neutral quality filter: applying global k-th percentile thresholds to pre-computed ccnet\_perplexity scores from RedPajama-V2 retains 86.4\% of Spanish documents but only 16.3\% of English documents at k=30 --- a 72.7 percentage-point gap driven by structural incomparability between CCNet's per-language KenLM perplexity scales. Across five threshold levels (k $\in$ {10, 20, 30, 40, 50}) on 208,262 documents spanning five European languages, we measure Cramér's V = 0.40--0.57 (all Holm-corrected p $\approx$ 0), with the Romance--Germanic language family ordering (es > fr > it > en > de) consistent across every threshold tested. We further find that literature-derived estimates of this bias underestimate the actual effect by 25--40\%, meaning correction studies sized from prior descriptions will be systematically underpowered. Our calibrated baseline provides the reference measurement needed for multilingual dataset curation research to evaluate quality-filter correction strategies on equal footing.

